% GENERATED FILE. Edit canonical lessons, then run npm run generate:books.
% canonical-source-hash: fnv1a64:97248b576acd93df
% canonical-lessons: TE-C248-molakettu, TE-C248-pituku, TE-C248-mosam-ceyu, TE-C248-docuko, TE-C248-appagincu, TE-R248-first-pass-words-from-chapters-240-243, TE-R248-second-pass-words-from-chapters-244-248

\chapter{To sprout, To milk, To cheat, To rob, To hand over}
\label{ch:te-to-sprout-to-milk-to-cheat-to-rob-to-hand-over}
\hlchaptermodality{248}

\begin{chapteropening}
\textbf{By the end of this chapter:} \emph{I can say to sprout, to milk, to cheat, to rob and to hand over.}

Complete the last lesson of chapter 248: I can say to sprout, to milk, to cheat, to rob and to hand over.
\end{chapteropening}

\section[\texorpdfstring{\te{మొలకెత్తు} (molakettu)}{మొలకెత్తు (molakettu)}]{\te{మొలకెత్తు} (molakettu) — to sprout}

\label{lesson:TE-C248-molakettu}



\begin{quote}
Before the new one: say the Telugu for to persuade, then the Telugu for to plead.
\end{quote}

\subsection*{You'll want to know: \te{మొలకెత్తు}}
\textbf{\te{మొలకెత్తు}} — \emph{molakettu} — \textquotedblleft{}to sprout\textquotedblright{}.

\begin{grammarlens}[title={What you've built}]
One more word for this chapter.
\end{grammarlens}

\subsection*{Your turn}
\begin{itemize}
\raggedright
  \item \emph{Say it:} \emph{molakettu}
  \item \emph{Say it:} \emph{molakettu}, once more
  \item \emph{Say it:} say \emph{molakettu}
  \item \emph{Recall:} say the Telugu for to persuade, then the Telugu for to plead, then say \emph{molakettu} again
\end{itemize}

\begin{tcolorbox}[breakable,colback=teal!4,colframe=teal!35!black,title={Before you move on}]
What does \te{మొలకెత్తు} mean? (\textquotedblleft{}To sprout\textquotedblright{}.) Say it once more.
\end{tcolorbox}
\section[\texorpdfstring{\te{పితుకు} (pituku)}{పితుకు (pituku)}]{\te{పితుకు} (pituku) — to milk (a cow)}

\label{lesson:TE-C248-pituku}



\begin{quote}
Before the new one: say the Telugu for to plead, then the Telugu for to sprout.
\end{quote}

\subsection*{You'll want to know: \te{పితుకు}}
\textbf{\te{పితుకు}} — \emph{pituku} — \textquotedblleft{}to milk (a cow)\textquotedblright{}.

\begin{grammarlens}[title={What you've built}]
One more word for this chapter.
\end{grammarlens}

\subsection*{Your turn}
\begin{itemize}
\raggedright
  \item \emph{Say it:} \emph{pituku}
  \item \emph{Say it:} \emph{pituku}, once more
  \item \emph{Say it:} say \emph{pituku}
  \item \emph{Recall:} say the Telugu for to plead, then the Telugu for to sprout, then say \emph{pituku} again
\end{itemize}

\begin{tcolorbox}[breakable,colback=teal!4,colframe=teal!35!black,title={Before you move on}]
What does \te{పితుకు} mean? (\textquotedblleft{}To milk (a cow)\textquotedblright{}.) Say it once more.
\end{tcolorbox}
\section[\texorpdfstring{\te{మోసం} \te{చేయు} (mōsaṁ cēyu)}{మోసం చేయు (mōsaṁ cēyu)}]{\te{మోసం} \te{చేయు} (mōsaṁ cēyu) — to cheat}

\label{lesson:TE-C248-mosam-ceyu}



\begin{quote}
Before the new one: say the Telugu for to sprout, then the Telugu for to milk.
\end{quote}

\subsection*{You'll want to know: \te{మోసం} \te{చేయు}}
\textbf{\te{మోసం} \te{చేయు}} — \emph{mōsaṁ cēyu} — \textquotedblleft{}to cheat\textquotedblright{}.

\begin{grammarlens}[title={What you've built}]
One more word for this chapter.
\end{grammarlens}

\subsection*{Your turn}
\begin{itemize}
\raggedright
  \item \emph{Say it:} \emph{mōsaṁ cēyu}
  \item \emph{Say it:} \emph{mōsaṁ cēyu}, once more
  \item \emph{Say it:} say \emph{mōsaṁ cēyu}
  \item \emph{Recall:} say the Telugu for to sprout, then the Telugu for to milk, then say \emph{mōsaṁ cēyu} again
\end{itemize}

\begin{tcolorbox}[breakable,colback=teal!4,colframe=teal!35!black,title={Before you move on}]
What does \te{మోసం} \te{చేయు} mean? (\textquotedblleft{}To cheat\textquotedblright{}.) Say it once more.
\end{tcolorbox}
\section[\texorpdfstring{\te{దోచుకో} (dōcukō)}{దోచుకో (dōcukō)}]{\te{దోచుకో} (dōcukō) — to rob}

\label{lesson:TE-C248-docuko}



\begin{quote}
Before the new one: say the Telugu for to milk, then the Telugu for to cheat.
\end{quote}

\subsection*{You'll want to know: \te{దోచుకో}}
\textbf{\te{దోచుకో}} — \emph{dōcukō} — \textquotedblleft{}to rob\textquotedblright{}.

\begin{grammarlens}[title={What you've built}]
One more word for this chapter.
\end{grammarlens}

\subsection*{Your turn}
\begin{itemize}
\raggedright
  \item \emph{Say it:} \emph{dōcukō}
  \item \emph{Say it:} \emph{dōcukō}, once more
  \item \emph{Say it:} say \emph{dōcukō}
  \item \emph{Recall:} say the Telugu for to milk, then the Telugu for to cheat, then say \emph{dōcukō} again
\end{itemize}

\begin{tcolorbox}[breakable,colback=teal!4,colframe=teal!35!black,title={Before you move on}]
What does \te{దోచుకో} mean? (\textquotedblleft{}To rob\textquotedblright{}.) Say it once more.
\end{tcolorbox}
\section[\texorpdfstring{\te{అప్పగించు} (appagiñcu)}{అప్పగించు (appagiñcu)}]{\te{అప్పగించు} (appagiñcu) — to hand over, to entrust}

\label{lesson:TE-C248-appagincu}



\begin{quote}
Before the new one: say the Telugu for to cheat, then the Telugu for to rob.
\end{quote}

\subsection*{You'll want to know: \te{అప్పగించు}}
\textbf{\te{అప్పగించు}} — \emph{appagiñcu} — \textquotedblleft{}to hand over, to entrust\textquotedblright{}.

\begin{grammarlens}[title={What you've built}]
One more word for this chapter.
\end{grammarlens}

\subsection*{Your turn}
\begin{itemize}
\raggedright
  \item \emph{Say it:} \emph{appagiñcu}
  \item \emph{Say it:} \emph{appagiñcu}, once more
  \item \emph{Say it:} say \emph{appagiñcu}
  \item \emph{Recall:} say the Telugu for to cheat, then the Telugu for to rob, then say \emph{appagiñcu} again
\end{itemize}

\begin{tcolorbox}[breakable,colback=teal!4,colframe=teal!35!black,title={Before you move on}]
What does \te{అప్పగించు} mean? (\textquotedblleft{}To hand over, to entrust\textquotedblright{}.) Say it once more.
\end{tcolorbox}
\section[(dialogue)]{Words from chapters 240-243 — a first pass}

\label{lesson:TE-R248-first-pass-words-from-chapters-240-243}



\begin{quote}
Say the words for to rob and to hand over.
\end{quote}

\subsection*{Your turn}
Say these aloud:

\begin{itemize}
\raggedright
  \item to lift, to carry, to pull, to push, to throw
  \item to swallow, to bite, to breathe in, to sneeze, to vomit
  \item to bring, to use, to finish, to begin, to iron
  \item to solve, to reduce, to raise, to inform, to transfer
\end{itemize}

\begin{tcolorbox}[breakable,colback=teal!4,colframe=teal!35!black,title={Before you move on}]
To lift? (\textbf{\te{ఎత్తు}}, \emph{ettu}.) To carry? (\textbf{\te{మోయు}}, \emph{mōyu}.) To pull? (\textbf{\te{లాగు}}, \emph{lāgu}.) To push? (\textbf{\te{నెట్టు}}, \emph{neṭṭu}.) To throw? (\textbf{\te{విసురు}}, \emph{visuru}.) To swallow? (\textbf{\te{మింగు}}, \emph{miṅgu}.) To bite? (\textbf{\te{కొరుకు}}, \emph{koruku}.) To breathe in? (\textbf{\te{పీల్చు}}, \emph{pīlcu}.) To sneeze? (\textbf{\te{తుమ్ము}}, \emph{tummu}.) To vomit? (\textbf{\te{వాంతి} \te{చేసుకో}}, \emph{vānti cēsukō}.) To bring? (\textbf{\te{తెచ్చు}}, \emph{teccu}.) To use? (\textbf{\te{ఉపయోగించు}}, \emph{upayōgiñcu}.) To finish? (\textbf{\te{ముగించు}}, \emph{mugiñcu}.) To begin? (\textbf{\te{మొదలుపెట్టు}}, \emph{modalupeṭṭu}.) To iron? (\textbf{\te{ఇస్త్రీ} \te{చేయు}}, \emph{istrī cēyu}.) To solve? (\textbf{\te{పరిష్కరించు}}, \emph{pariṣkariñcu}.) To reduce? (\textbf{\te{తగ్గించు}}, \emph{taggiñcu}.) To raise? (\textbf{\te{పెంచు}}, \emph{peñcu}.) To inform? (\textbf{\te{తెలియజేయు}}, \emph{teliyajēyu}.) To transfer? (\textbf{\te{బదిలీ} \te{చేయు}}, \emph{badilī cēyu}.)
\end{tcolorbox}
\section[(dialogue)]{Words from chapters 244-248 — a second pass}

\label{lesson:TE-R248-second-pass-words-from-chapters-244-248}



\begin{quote}
Say the words for to blow and to shine.
\end{quote}

\subsection*{Your turn}
Say these aloud:

\begin{itemize}
\raggedright
  \item to blow, to shine, to freeze, to wilt, to bark
  \item to collect, to broadcast, to print, to download, to type
  \item to lend, to give back, to empty, to decrease, to refuse
  \item to pinch, to scratch, to whisper, to persuade, to plead
  \item to sprout, to milk, to cheat, to rob, to hand over
\end{itemize}

\begin{tcolorbox}[breakable,colback=teal!4,colframe=teal!35!black,title={Before you move on}]
To blow? (\textbf{\te{వీచు}}, \emph{vīcu}.) To shine? (\textbf{\te{ప్రకాశించు}}, \emph{prakāśiñcu}.) To freeze? (\textbf{\te{గడ్డకట్టు}}, \emph{gaḍḍakaṭṭu}.) To wilt? (\textbf{\te{వాడిపోవు}}, \emph{vāḍipōvu}.) To bark? (\textbf{\te{మొరుగు}}, \emph{morugu}.) To collect? (\textbf{\te{సేకరించు}}, \emph{sēkariñcu}.) To broadcast? (\textbf{\te{ప్రసారం} \te{చేయు}}, \emph{prasāraṁ cēyu}.) To print? (\textbf{\te{ముద్రించు}}, \emph{mudriñcu}.) To download? (\textbf{\te{డౌన్లోడ్} \te{చేయు}}, \emph{ḍaunlōḍ cēyu}.) To type? (\textbf{\te{టైప్} \te{చేయు}}, \emph{ṭaip cēyu}.) To lend? (\textbf{\te{అప్పు} \te{ఇవ్వు}}, \emph{appu ivvu}.) To give back? (\textbf{\te{తిరిగి} \te{ఇవ్వు}}, \emph{tirigi ivvu}.) To empty? (\textbf{\te{ఖాళీ} \te{చేయు}}, \emph{khāḷī cēyu}.) To decrease? (\textbf{\te{తగ్గు}}, \emph{taggu}.) To refuse? (\textbf{\te{నిరాకరించు}}, \emph{nirākariñcu}.) To pinch? (\textbf{\te{గిల్లు}}, \emph{gillu}.) To scratch? (\textbf{\te{గోకు}}, \emph{gōku}.) To whisper? (\textbf{\te{గుసగుసలాడు}}, \emph{gusagusalāḍu}.) To persuade? (\textbf{\te{ఒప్పించు}}, \emph{oppiñcu}.) To plead? (\textbf{\te{బతిమాలు}}, \emph{batimālu}.) To sprout? (\textbf{\te{మొలకెత్తు}}, \emph{molakettu}.) To milk? (\textbf{\te{పితుకు}}, \emph{pituku}.) To cheat? (\textbf{\te{మోసం} \te{చేయు}}, \emph{mōsaṁ cēyu}.) To rob? (\textbf{\te{దోచుకో}}, \emph{dōcukō}.) To hand over? (\textbf{\te{అప్పగించు}}, \emph{appagiñcu}.)
\end{tcolorbox}
